\documentclass[12pt,a4paper]{article}

\usepackage[spanish,es-nodecimaldot]{babel}
\usepackage[utf8]{inputenc}
\usepackage[T1]{fontenc}
\usepackage{lmodern}
\usepackage{amsmath,amssymb}
\usepackage{graphicx}
\usepackage{booktabs}
\usepackage{hyperref}
\usepackage{enumitem}
\usepackage{geometry}
\usepackage{listings}
\usepackage{xcolor}

\geometry{margin=2.5cm}

\lstset{
  basicstyle=\footnotesize\ttfamily,
  breaklines=true,
  frame=single,
}

\begin{document}
\section{04 --- Validación y lectura correcta de resultados}

\subsection{1) Por qué no basta con una accuracy alta}

Un resultado alto en split aleatorio puede ocultar problemas de generalización. Por eso el proyecto
incluye varias validaciones con funciones distintas.

\subsection{2) Check A (evaluación directa)}

Archivo: \verb|src/validate_checks.py| $\to$ \verb|check_a_direct_eval|

\begin{itemize}
  \item compara \verb|model.predict(X_test[i])| con \verb|y_test[i]|
  \item evita depender de mecánicas internas del entorno para computar métrica
\end{itemize}

Resultado histórico (artefacto comprometido):

\begin{itemize}
  \item accuracy $\approx$ \texttt{0.9939}
\end{itemize}

\subsection{3) Check B (anti-leakage con etiquetas barajadas)}

Archivo: \verb|src/validate_checks.py| $\to$ \verb|check_b_shuffled_labels|

\begin{itemize}
  \item baraja \texttt{y\_train}
  \item reentrena brevemente
  \item espera rendimiento cercano a baseline de clase mayoritaria
\end{itemize}

Resultado histórico comprometido:

\begin{itemize}
  \item \texttt{shuffled\_accuracy = 0.4773}
  \item \texttt{leakage\_detected = false}
\end{itemize}

\subsection{4) Check C (split duro por CSV/día)}

Archivo: \verb|src/validate_checks.py| $\to$ \verb|check_c_csv_split|

\begin{itemize}
  \item train en días/patrones distintos de test
  \item prueba de generalización más realista
\end{itemize}

Resultado histórico comprometido:

\begin{itemize}
  \item accuracy $\approx$ \texttt{0.84135}
  \item recall ataque y F1 bajan de forma notable frente a random split
\end{itemize}

Lectura correcta: el pipeline funciona muy bien in-distribution, pero la generalización dura es
sustancialmente más difícil.

\subsection{5) Leave-one-exact-CSV-out}

Archivo: \verb|src/validate_leave_one_csv_out.py|

\begin{itemize}
  \item holdout de un CSV exacto por fold
  \item agrega métricas por fold y globales
  \item incluye métricas operativas (\texttt{balanced\_accuracy}, \texttt{fpr}, \texttt{fnr},
        \texttt{block\_rate}, \texttt{reward\_per\_sample}, tiempos)
\end{itemize}

Estado a defender con rigor:

\begin{itemize}
  \item workflow implementado en código
  \item sin artefacto agregado completo comprometido aún
\end{itemize}

\subsection{6) Cómo explicar contradicciones aparentes en métricas}

No hay contradicción entre ``99.8\% en random split'' y ``84\% en split duro''. Son preguntas
distintas:

\begin{itemize}
  \item random split: rendimiento en distribución muy parecida
  \item split duro: robustez a cambios de día/patrón
\end{itemize}

La defensa sólida es mostrar ambas, no esconder la más dura.

\subsection{7) Regla metodológica para el tribunal}

Siempre separar:

\begin{enumerate}
  \item defaults actuales del código
  \item configuración real del run histórico citado
  \item resultado medido respaldado por artefacto
\end{enumerate}

\end{document}
